\documentclass[a4paper,12pt,twocolumn]{article}

\usepackage[landscape,top=1cm,bottom=1.4cm,left=1cm,right=1cm,columnsep=1cm]{geometry}
\usepackage[fleqn]{amsmath}
\usepackage{amssymb}
\setlength{\mathindent}{1.5em}
\usepackage{fontspec}
\usepackage{enumitem}
\usepackage{framed}
\usepackage{needspace}
\usepackage{xeCJK}

% Set fonts
\setmainfont{Times New Roman}
\setsansfont{Arial}
\setmonofont{JetBrains Mono}

% Chinese font support
\setCJKmainfont{Iansui}[
  BoldFont=Iansui,
  AutoFakeBold=3.17
]

\pagestyle{plain}
\setlength{\columnseprule}{0.2pt}
\newcommand{\examdate}{2026/03/11}
\newenvironment{solutionbox}{\begin{framed}}{\end{framed}}

\begin{document}

\twocolumn[
{\centering
\Large\bfseries Calculus Problems Solutions\par
\vspace{0.5em}
\noindent\hfill\normalsize\normalfont \examdate\par
\vspace{0.8em}
}
]

\begin{enumerate}[label=\textbf{\arabic*.}, leftmargin=2em, itemsep=1.2em]
\setcounter{enumi}{48}

\Needspace{28\baselineskip}
\item Determine whether the series is convergent or divergent. If it is convergent, find its sum:
\[
\sum_{n=1}^{\infty}\left(\frac{1}{e^n}+\frac{1}{n(n+1)}\right)
\]

\begin{solutionbox}
\textbf{Method 1: Split the series into two known series}

Write
\[
\sum_{n=1}^{\infty}\left(\frac{1}{e^n}+\frac{1}{n(n+1)}\right)
=
\sum_{n=1}^{\infty}\frac{1}{e^n}
+
\sum_{n=1}^{\infty}\frac{1}{n(n+1)}.
\]
For the first series,
\[
\sum_{n=1}^{\infty}\frac{1}{e^n}
=
\sum_{n=1}^{\infty}\left(\frac{1}{e}\right)^n
=
\frac{1/e}{1-1/e}
=
\frac{1}{e-1}.
\]
For the second series, use partial fractions:
\[
\frac{1}{n(n+1)}=\frac{1}{n}-\frac{1}{n+1}.
\]
So
\[
\sum_{n=1}^{\infty}\frac{1}{n(n+1)}
=
\sum_{n=1}^{\infty}\left(\frac{1}{n}-\frac{1}{n+1}\right).
\]
This is a telescoping series. Its \(N\)-th partial sum is
\[
S_N=
\left(1-\frac12\right)+\left(\frac12-\frac13\right)+\cdots+\left(\frac1N-\frac1{N+1}\right)
=
1-\frac{1}{N+1}.
\]
Taking \(N\to\infty\), we get
\[
\sum_{n=1}^{\infty}\frac{1}{n(n+1)}=1.
\]
Therefore,
\[
\sum_{n=1}^{\infty}\left(\frac{1}{e^n}+\frac{1}{n(n+1)}\right)
=
\frac{1}{e-1}+1
=
\frac{e}{e-1}.
\]
\end{solutionbox}

\begin{solutionbox}
\textbf{Method 2: Work with partial sums directly}

Let
\[
S_N=\sum_{n=1}^{N}\left(\frac{1}{e^n}+\frac{1}{n(n+1)}\right).
\]
Then
\[
S_N=
\sum_{n=1}^{N}\frac{1}{e^n}
+
\sum_{n=1}^{N}\frac{1}{n(n+1)}.
\]
For the finite geometric sum,
\[
\sum_{n=1}^{N}\frac{1}{e^n}
=
\sum_{n=1}^{N}\left(\frac{1}{e}\right)^n
=
\frac{\frac1e\left(1-\left(\frac1e\right)^N\right)}{1-\frac1e}
=
\frac{1-(1/e)^N}{e-1}.
\]
For the second part,
\[
\frac{1}{n(n+1)}=\frac{1}{n}-\frac{1}{n+1},
\]
so
\[
\sum_{n=1}^{N}\frac{1}{n(n+1)}
=
\sum_{n=1}^{N}\left(\frac{1}{n}-\frac{1}{n+1}\right)
=
1-\frac{1}{N+1}.
\]
Hence
\[
S_N=
\frac{1-(1/e)^N}{e-1}+1-\frac{1}{N+1}.
\]
Now let \(N\to\infty\). Since
\[
\left(\frac{1}{e}\right)^N\to 0
\qquad\text{and}\qquad
\frac{1}{N+1}\to 0,
\]
we obtain
\[
\lim_{N\to\infty}S_N
=
\frac{1}{e-1}+1
=
\frac{e}{e-1}.
\]

\medskip
\textbf{Answer:}
\[
\boxed{\sum_{n=1}^{\infty}\left(\frac{1}{e^n}+\frac{1}{n(n+1)}\right)\text{ converges, and its sum is }\frac{e}{e-1}.}
\]
\end{solutionbox}

\Needspace{20\baselineskip}
\item Determine whether the series is convergent or divergent. If it is convergent, find its sum:
\[
\sum_{n=1}^{\infty}\frac{e^n}{n^2}
\]

\begin{solutionbox}
\textbf{Method 1: nth-term test with L'Hôpital's Rule}

Let
\[
a_n=\frac{e^n}{n^2}.
\]
A necessary condition for the convergence of \(\sum a_n\) is
\[
\lim_{n\to\infty} a_n = 0.
\]
So we examine the limit of \(a_n\).

Consider the reciprocal:
\[
\lim_{n\to\infty}\frac{n^2}{e^n}
=
\lim_{x\to\infty}\frac{x^2}{e^x}.
\]
This is an \(\infty/\infty\) form, so apply L'Hôpital's Rule twice:
\[
\lim_{x\to\infty}\frac{x^2}{e^x}
=
\lim_{x\to\infty}\frac{2x}{e^x}
=
\lim_{x\to\infty}\frac{2}{e^x}
=0.
\]
Therefore,
\[
\lim_{n\to\infty}\frac{n^2}{e^n}=0
\quad\Longrightarrow\quad
\lim_{n\to\infty}\frac{e^n}{n^2}=\infty.
\]
Hence \(a_n \not\to 0\). Since the terms of the series do not approach \(0\), the series diverges.
\end{solutionbox}

\begin{solutionbox}
\textbf{Method 2: Exponential growth dominates polynomial growth}

Let
\[
a_n=\frac{e^n}{n^2}.
\]
The numerator \(e^n\) has exponential growth, while the denominator \(n^2\) has polynomial growth. Exponential growth is much faster than polynomial growth, so
\[
\frac{e^n}{n^2}\to\infty
\qquad (n\to\infty).
\]
In particular,
\[
\lim_{n\to\infty}\frac{e^n}{n^2}\neq 0.
\]
Therefore the necessary condition for convergence fails, so the series diverges.

\end{solutionbox}

\setcounter{enumi}{76}
\Needspace{24\baselineskip}
\item Find the value of \(c\) if
\[
\sum_{n=2}^{\infty}(1+c)^{-n}=2
\]

\begin{solutionbox}
\textbf{Method 1: Treat it as a geometric series directly}

Let
\[
x=1+c.
\]
Then
\[
\sum_{n=2}^{\infty}(1+c)^{-n}
=
\sum_{n=2}^{\infty}x^{-n}
\]
is a geometric series with first term \(a=x^{-2}\) and common ratio \(r=x^{-1}\). Hence
\[
\frac{x^{-2}}{1-x^{-1}}=2.
\]
Simplify:
\[
\frac{1/x^2}{1-1/x}
=
\frac{1}{x(x-1)}
=2.
\]
So
\[
1=2x(x-1)
\quad\Longrightarrow\quad
2x^2-2x-1=0.
\]
Using the quadratic formula,
\[
x=\frac{2\pm\sqrt{(-2)^2-4(2)(-1)}}{4}
=
\frac{1\pm\sqrt{3}}{2}.
\]
For the infinite geometric series to converge, we need
\[
|r|<1
\quad\Longleftrightarrow\quad
\left|\frac{1}{x}\right|<1
\quad\Longleftrightarrow\quad
|x|>1.
\]
Thus the valid value is
\[
x=\frac{1+\sqrt{3}}{2}.
\]
Therefore,
\[
c=x-1=\frac{1+\sqrt{3}}{2}-1=\frac{\sqrt{3}-1}{2}.
\]
\end{solutionbox}

\begin{solutionbox}
\textbf{Method 2: Substitute \(r=(1+c)^{-1}\)}

Let
\[
r=(1+c)^{-1}.
\]
Then
\[
\sum_{n=2}^{\infty}(1+c)^{-n}
=
\sum_{n=2}^{\infty}r^n
=
\frac{r^2}{1-r}=2,
\]
with convergence condition \(|r|<1\).

So
\[
r^2=2(1-r)
\quad\Longrightarrow\quad
r^2+2r-2=0.
\]
Hence
\[
r=\frac{-2\pm\sqrt{4+8}}{2}
=
-1\pm\sqrt{3}.
\]
Since \(|r|<1\), the valid value is
\[
r=\sqrt{3}-1.
\]
Now
\[
1+c=\frac{1}{r}=\frac{1}{\sqrt{3}-1}
=
\frac{\sqrt{3}+1}{2},
\]
so
\[
c=\frac{\sqrt{3}+1}{2}-1=\frac{\sqrt{3}-1}{2}.
\]

\medskip
\textbf{Answer:}
\[
\boxed{c=\frac{\sqrt{3}-1}{2}}
\]
\end{solutionbox}

\setcounter{enumi}{7}
\Needspace{28\baselineskip}
\item Use the Integral Test to determine whether the series is convergent or divergent:
\[
\sum_{n=1}^{\infty} n^2 e^{-n^3}
\]

\begin{solutionbox}
\textbf{Method 1: Integral Test}

Let
\[
f(x)=x^2e^{-x^3}, \qquad x\ge 1.
\]
For \(x\ge 1\), \(f(x)>0\) and \(f(x)\) is continuous. To use the Integral Test, we also check that \(f(x)\) is decreasing.

Differentiate:
\[
f'(x)=\frac{d}{dx}\left(x^2e^{-x^3}\right)
=2xe^{-x^3}+x^2e^{-x^3}(-3x^2).
\]
So
\[
f'(x)=e^{-x^3}(2x-3x^4)=xe^{-x^3}(2-3x^3).
\]
For \(x\ge 1\), we have \(2-3x^3<0\), while \(x>0\) and \(e^{-x^3}>0\). Hence
\[
f'(x)<0 \qquad \text{for } x\ge 1,
\]
so \(f(x)\) is decreasing.

Now evaluate the improper integral:
\[
\int_1^\infty x^2e^{-x^3}\,dx.
\]
Use the substitution \(u=x^3\), so \(du=3x^2\,dx\) and \(x^2\,dx=\frac13\,du\). Then
\[
\int_1^\infty x^2e^{-x^3}\,dx
=
\frac13\int_1^\infty e^{-u}\,du
=
\frac13\left[-e^{-u}\right]_1^\infty
=
\frac{1}{3e}.
\]
Since this improper integral is finite, the Integral Test shows that
\[
\sum_{n=1}^{\infty} n^2 e^{-n^3}
\]
converges.
\end{solutionbox}

\begin{solutionbox}
\textbf{Method 2: Comparison with a convergent geometric series}

For sufficiently large \(n\), we have
\[
n^2\le e^n.
\]
Therefore,
\[
n^2e^{-n^3}\le e^n e^{-n^3}=e^{\,n-n^3}.
\]
Since \(n^3-n\ge n\) for all sufficiently large \(n\),
\[
e^{\,n-n^3}\le e^{-n}.
\]
Hence, for all sufficiently large \(n\),
\[
0\le n^2e^{-n^3}\le e^{-n}.
\]
But
\[
\sum_{n=1}^{\infty} e^{-n}
\]
is a geometric series with ratio \(1/e<1\), so it converges. By the Comparison Test,
\[
\sum_{n=1}^{\infty} n^2 e^{-n^3}
\]
also converges.

\medskip
\textbf{Answer:}
\[
\boxed{\sum_{n=1}^{\infty} n^2 e^{-n^3}\text{ converges}.}
\]
\end{solutionbox}

\setcounter{enumi}{9}
\Needspace{30\baselineskip}
\item Use the Integral Test to determine whether the series is convergent or divergent:
\[
\sum_{n=1}^{\infty}\frac{\tan^{-1} n}{1+n^2}
\]

\begin{solutionbox}
\textbf{Method 1: Integral Test}

Let
\[
f(x)=\frac{\tan^{-1}x}{1+x^2}, \qquad x\ge 1.
\]
For \(x\ge 1\), \(f(x)>0\) and \(f(x)\) is continuous. To use the Integral Test, we also check that \(f(x)\) is decreasing.

Differentiate:
\[
f'(x)=\frac{(1+x^2)\cdot \frac{1}{1+x^2}-(\tan^{-1}x)(2x)}{(1+x^2)^2}
=\frac{1-2x\tan^{-1}x}{(1+x^2)^2}.
\]
For \(x\ge 1\),
\[
\tan^{-1}x\ge \tan^{-1}1=\frac{\pi}{4},
\]
so
\[
2x\tan^{-1}x\ge 2\cdot 1\cdot \frac{\pi}{4}=\frac{\pi}{2}>1.
\]
Hence
\[
1-2x\tan^{-1}x<0.
\]
Since the denominator is always positive,
\[
f'(x)<0 \qquad \text{for } x\ge 1,
\]
so \(f(x)\) is decreasing.

Now evaluate the improper integral:
\[
\int_1^\infty \frac{\tan^{-1}x}{1+x^2}\,dx.
\]
Use the substitution
\[
u=\tan^{-1}x, \qquad du=\frac{1}{1+x^2}\,dx.
\]
Then
\[
\int_1^\infty \frac{\tan^{-1}x}{1+x^2}\,dx
=
\int_{\pi/4}^{\pi/2} u\,du
=
\left[\frac{u^2}{2}\right]_{\pi/4}^{\pi/2}
=
\frac{3\pi^2}{32}.
\]
Since this improper integral is finite, the Integral Test shows that
\[
\sum_{n=1}^{\infty}\frac{\tan^{-1} n}{1+n^2}
\]
converges.
\end{solutionbox}

\begin{solutionbox}
\textbf{Method 2: Comparison Test}

For every \(n\ge 1\),
\[
\tan^{-1}n<\frac{\pi}{2}.
\]
Therefore,
\[
0\le \frac{\tan^{-1}n}{1+n^2}
\le \frac{\pi/2}{1+n^2}.
\]
Also, since \(1+n^2\ge n^2\),
\[
\frac{1}{1+n^2}\le \frac{1}{n^2}.
\]
Hence
\[
0\le \frac{\tan^{-1}n}{1+n^2}
\le \frac{\pi}{2n^2}.
\]
Now
\[
\sum_{n=1}^{\infty}\frac{\pi}{2n^2}
=
\frac{\pi}{2}\sum_{n=1}^{\infty}\frac{1}{n^2}
\]
converges because \(\sum 1/n^2\) is a \(p\)-series with \(p=2>1\). By the Comparison Test,
\[
\sum_{n=1}^{\infty}\frac{\tan^{-1} n}{1+n^2}
\]
also converges.

\medskip
\textbf{Answer:}
\[
\boxed{\sum_{n=1}^{\infty}\frac{\tan^{-1} n}{1+n^2}\text{ converges}.}
\]
\end{solutionbox}

\end{enumerate}

\end{document}
